\chapter{Sacroiliac Joint Pain}

\section*{Definition}
Sacroiliac (SI) joint pain is pain originating from the articulation between the sacrum and ilium, typically presenting as unilateral low back pain below L5, with referral to the buttock, groin, and posterior thigh. Diagnosis is confirmed by a positive response to a series of provocative maneuvers and diagnostic intra-articular injection.

\section*{Pathophysiology}
The SI joint is a diarthrodial joint with a synovial capsule, supported by strong anterior and interosseous ligaments. It undergoes approximately 2--4 degrees of rotation and translation during gait. Joint dysfunction may arise from hypomobility (fixation) or hypermobility (excessive laxity). Abnormal shear forces across the joint irritate the innervated posterior ligamentous complex and dorsal rami of S1--S4 nerve roots. Inflammatory sacroiliitis involves subchondral bone marrow edema, erosion, and ankylosis.

\section*{Causes}
\begin{itemize}
  \item \textbf{Inflammatory:} Axial spondyloarthropathy (ankylosing spondylitis, psoriatic arthritis, reactive arthritis, IBD-associated arthritis)
  \item \textbf{Traumatic:} Falls onto the buttocks, direct blow, motor vehicle accident
  \item \textbf{Pregnancy:} Hormone-induced ligamentous laxity, childbirth trauma
  \item \textbf{Iatrogenic:} Post-lumbar fusion (adjacent segment disease), post-hip arthroplasty
  \item \textbf{Degenerative:} Osteoarthritis, osteophyte formation, joint space narrowing
  \item \textbf{Infectious:} Septic sacroiliitis (rare, typically hematogenous spread)
  \item \textbf{Biomechanical:} Leg length discrepancy, abnormal gait pattern, prolonged asymmetric postures
\end{itemize}

\section*{When to See a Doctor}
See your doctor if:
\begin{itemize}
  \item Unilateral low back or buttock pain below L5
  \item Pain with sitting, rising from a chair, or rolling over in bed
  \item Positive provocative tests (FABER, Gaenslen, thigh thrust, compression/distraction)
  \item Morning stiffness lasting > 30 minutes with improvement after activity
\end{itemize}

\section*{Related Symptoms}
\begin{itemize}
  \item Unilateral buttock pain, often alternating sides in inflammatory disease
  \item Referral to the posterior thigh (not below the knee typically)
  \item Pain with transitional movements (sit-to-stand, stair climbing)
  \item Tenderness over the posterior superior iliac spine and SI joint line
  \item Antalgic gait with decreased stance phase on the affected side
\end{itemize}
\textbf{Important:} SI joint pain is a common mimic of lumbar radiculopathy and hip pathology; a positive cluster of 3 or more provocative tests has high diagnostic accuracy.

\section*{Self-Care / Home Management}
\begin{itemize}
  \item Activity modification to avoid painful movements (e.g., single-leg stance, twisting)
  \item Ice or heat packs applied over the SI joint region for 15--20 minutes
  \item Over-the-counter NSAIDs for pain and inflammation
  \item SI joint belt to provide external compression and stability
  \item Avoid sitting on uneven surfaces or with wallet in back pocket
\end{itemize}

\section*{How Your Doctor Will Evaluate You}
\begin{itemize}
  \item Provocative physical examination tests: FABER, Gaenslen, thigh thrust, compression, distraction
  \item Diagnostic SI joint injection (image-guided with local anesthetic) as the gold standard
  \item MRI with STIR sequences for bone marrow edema (active sacroiliitis)
  \item CT for evaluation of erosions, subchondral sclerosis, and ankylosis
  \item Plain radiography: Ferguson view for joint space narrowing and erosions
  \item Laboratory studies: ESR, CRP, HLA-B27 if inflammatory spondyloarthropathy is suspected
\end{itemize}

\section*{Treatment Options}
\begin{itemize}
  \item NSAIDs and physical therapy (core strengthening, pelvic stabilization, manual therapy) as first-line
  \item Image-guided intra-articular corticosteroid injection for diagnostic confirmation and therapy
  \item Radiofrequency ablation of the lateral branches (S1--S4) for chronic refractory pain
  \item Disease-modifying antirheumatic drugs (sulfasalazine, methotrexate) or TNF inhibitors for inflammatory sacroiliitis
  \item SI joint fusion for refractory cases after failure of non-surgical management
  \item Treatment of underlying leg length discrepancy with shoe lift
\end{itemize}

\clearpage
